El teclado, los cursores y el joystick 1 funcionan a la vez.

\begin{center}
\begin{tabular}{@{}lll@{}}
\toprule
\textbf{\tAction} & \textbf{\tKeyboard} & \textbf{\tJoystick}\\
\midrule
Girar a izquierda / derecha & \key{$\leftarrow$} \key{$\rightarrow$} & palanca izquierda / derecha\\
Acelerar & \key{$\uparrow$} & palanca arriba\\
Frenar & \key{$\downarrow$} & palanca abajo\\
Cambiar de marcha, \texttt{LO} / \texttt{HI} (\textbf{MARCHA}) & \key{ESPACIO} & botón 1\\
Turbo & \key{M} & botón 2\\
Pausa & \key{P} & \\
Abandonar la partida & \key{Q} & \\
Empezar, saltar el aviso, continuar (\textbf{MARCHA}) & \key{ESPACIO} & botón 1\\
\bottomrule
\end{tabular}
\end{center}

\textbf{Notas.}
\begin{itemize}
\item Cada pulsación de \textbf{MARCHA} cambia entre la marcha corta, \texttt{LO}, y la larga, \texttt{HI}. Arranca en \texttt{LO} y cambia a \texttt{HI} cuando la velocidad deje de subir.
\item El \textbf{turbo} da un acelerón durante unos segundos. Tienes \textbf{tres por fase}; los iconos de abajo a la izquierda muestran cuántos quedan.
\item \key{P} pone el juego en pausa y lo deja en silencio. Pulsa cualquier control para seguir.
\item \key{Q} acaba la partida en el acto, como si se agotara el tiempo sin créditos.
\item Pulsar a la vez izquierda y derecha, o acelerar y frenar, no hace nada.
\end{itemize}
